% ===========================================================
%  THEOREM 2 -- FORMAL SETUP
%  v15 NeurIPS 2026 / ICML 2027 paper
%  Author: Seungho Cook
% ===========================================================
\documentclass[11pt,a4paper]{article}
\usepackage[margin=1in]{geometry}
\usepackage{amsmath,amssymb,amsthm,bm,hyperref,booktabs}
\newtheorem{theorem}{Theorem}
\newtheorem{proposition}{Proposition}
\newtheorem{corollary}{Corollary}
\newtheorem{lemma}{Lemma}
\newtheorem{definition}{Definition}
\newcommand{\R}{\mathbb{R}}
\newcommand{\E}{\mathbb{E}}
\newcommand{\KL}{\mathrm{KL}}
\newcommand{\TODO}[1]{\textbf{\textcolor{red}{[TODO: Seungho verify -- #1]}}}
\usepackage{xcolor}
\title{Theorem~2 Setup: Unified Shift Decomposition under Linear Correction\\
\large v15 NeurIPS/ICML chapter, 2026-04-25}
\author{Seungho Cook}
\begin{document}
\maketitle

\section{Joint Distributions and Correction Operator}

\subsection{Measurable spaces}

Let $\mathcal{X}\subseteq\R^{p}$ be a feature space,
$\mathcal{Y}=\{0,1\}$ a binary label space (the multi-class extension is
deferred to Appendix~D), and $\mathcal{B}=\{1,\dots,K\}$ a finite set of
\emph{batch/domain identifiers}. A sample is a triple
$(X,Y,B)\in\mathcal{X}\times\mathcal{Y}\times\mathcal{B}$.

\subsection{Source and target distributions}

We consider two probability measures on
$\mathcal{X}\times\mathcal{Y}\times\mathcal{B}$:
\[
  p_s(x,y,b) \;=\; p_s(x\mid y,b)\,p_s(y\mid b)\,p_s(b),
  \qquad
  p_t(x,y,b) \;=\; p_t(x\mid y,b)\,p_t(y\mid b)\,p_t(b).
\]
The marginals
$p_s(x)=\sum_{y,b}p_s(x,y,b)$ and $p_t(x)$ are the \emph{observable
covariate distributions}; $p_s(y),p_t(y)$ are the \emph{label
distributions}; and $p_s(x\mid y), p_t(x\mid y)$ are the
\emph{class-conditional distributions}, obtained by marginalising out
$B$.

\subsection{Correction operator}

\begin{definition}[Linear correction operator]
\label{def:linear-correction}
A \emph{linear correction operator} is a batch-wise affine map
\[
  T: \mathcal{X}\times\mathcal{B} \to \mathcal{X},
  \qquad
  T(x,b) \;=\; A_b\,x + c_b,
\]
where $A_b\in\R^{p\times p}$ is non-singular and $c_b\in\R^{p}$ for
each batch $b\in\mathcal{B}$. The image under $T$ of a random variable
$(X,B)$ is denoted $\tilde X = T(X,B)$.
\end{definition}

\paragraph{Examples.}
\begin{itemize}
\item \textbf{ComBat (Johnson \emph{et al.}\ 2007).} With
  $A_b = \mathrm{diag}(1/\sqrt{\gamma_b})$ and
  $c_b = -\mu_b/\sqrt{\gamma_b}$, one recovers location-and-scale
  ComBat (empirical Bayes priors yield slightly shrunken estimates; the
  result remains affine per batch).
\item \textbf{CORAL (Sun \emph{et al.}\ 2016).} A single global
  whitening matrix applied to the target, $T(x,b)=A x+c$, independent
  of $b$.
\item \textbf{Mean centring.} $A_b=I$, $c_b=-\mu_b$.
\item \textbf{Subspace projection.} $T(x,b)=(I-P_{\mathcal{B}})(x-\mu_b)$
  where $P_{\mathcal{B}}$ projects onto a batch subspace
  $\mathcal{B}\subset\R^{p}$; this is the Leek--Storey SVA (2007) linear
  surrogate-variable approach.
\end{itemize}

\subsection{Label and batch directions}

Let $w_Y\in\R^{p}$ solve the population Fisher discriminant between
classes under $p_s$:
\[
  w_Y \;=\; \Sigma_s^{-1}\bigl(\mu_{s,1}-\mu_{s,0}\bigr),
  \qquad
  \mu_{s,y}=\E_{p_s}[X\mid Y=y],\;
  \Sigma_s=\mathrm{Cov}_{p_s}(X).
\]
Let $w_B\in\R^{p}$ be the top right-singular vector of the
batch-centroid matrix $M=[\mu_{s,b}-\bar\mu_s]_{b\in\mathcal{B}}$, so
$w_B$ is the \emph{dominant batch direction}.

\subsection{Alignment angle}

Define the (population) \emph{alignment coefficient}
\[
  \alpha \;=\; \frac{\langle w_Y, w_B\rangle^{2}}
                    {\|w_Y\|^{2}\,\|w_B\|^{2}} \;\in\;[0,1].
\]
The event $\alpha>\tfrac12$ is Theorem~1's flip condition (v10 chapter
and AAAI preprint); under $\alpha=0$ the label and batch axes are
orthogonal and classical ComBat is consistent.

\section{DIAL Metric (recap)}

\begin{definition}[DIAL]
\label{def:dial}
Let $\mathrm{AUC}(\tilde X,Y)$ denote cross-validated AUC for a fixed
classifier family $\mathcal{F}$ trained on $\tilde X$ to predict $Y$,
and let $\mathrm{AUC}^{\flip}=\max(\mathrm{AUC},1-\mathrm{AUC})$. The
\emph{direction-invariant AUC leakage} is
\[
  \mathrm{DIAL}(\tilde X,Y) \;=\; \bigl(\mathrm{AUC}^{\flip}
  (\tilde X,Y)-\tfrac12\bigr)\,\mathbf{1}\!\left[\mathrm{AUC}
  (\tilde X,Y)<\tfrac12\right].
\]
\end{definition}

Equivalently, DIAL is strictly positive iff the classifier trained on
corrected features achieves AUC $<\tfrac12$, i.e.\ the posterior
direction has \emph{flipped} relative to the true label (``leakage
through the label axis'').

\section{Three shift components}

Given $p_s,p_t$ and a correction $T$, we define three Kullback--Leibler
shift components on the corrected variable $\tilde X=T(X,B)$:
\begin{align}
  \Delta_{\text{cov}} &:= \KL\!\bigl(p_s(\tilde x)\,\|\,p_t(\tilde x)\bigr),
    \label{eq:cov}\\[2pt]
  \Delta_{\text{lab}} &:= \KL\!\bigl(p_s(y)\,\|\,p_t(y)\bigr),
    \label{eq:lab}\\[2pt]
  \Delta_{\text{cond}} &:= \E_{y\sim p_s}\bigl[\KL\!\bigl(
      p_s(\tilde x\mid y)\,\|\,p_t(\tilde x\mid y)\bigr)\bigr].
    \label{eq:cond}
\end{align}

\paragraph{Ben-David / Quionero-Candela taxonomy.}
\[
  \text{Shift type} \;=\;
    \begin{cases}
      \text{covariate}   & \Delta_{\text{cov}}>0,\; \Delta_{\text{lab}}=\Delta_{\text{cond}}=0,\\
      \text{label (prior)}     & \Delta_{\text{lab}}>0,\; \Delta_{\text{cond}}=0,\\
      \text{concept}     & p_s(y\mid \tilde x)\neq p_t(y\mid \tilde x),\\
      \text{conditional} & \Delta_{\text{cond}}>0 \text{ (subspace-aligned: ours).}
    \end{cases}
\]
The ``subspace-aligned conditional'' case is the one not covered by
Shimodaira's classical covariate-shift theory or by pure label shift,
and is the topic of Theorem~2 below.

\section{What Theorem 2 will establish}

The chapter's main theoretical claim, stated formally in
\texttt{theorem2\_proof.tex}, is:
\begin{itemize}
\item DIAL$(\tilde X,Y)>0$ is governed by the component of
  $\Delta_{\text{cond}}$ that lies \emph{parallel} to the span of the
  label direction $w_Y$ after correction;
\item when correction $T$ acts primarily \emph{inside}
  $\mathrm{span}(w_Y)$, the flip of Theorem~1 is exactly the event
  ``DIAL $>$ 0'';
\item DIAL therefore \emph{unifies} the three Ben-David types into a
  single decomposition and fires \emph{only} on type-(d), the
  subspace-aligned conditional shift.
\end{itemize}

\paragraph{Honest scope.}
\TODO{This setup assumes (i) Gaussian class-conditionals for the
cleanest KL chain rule, (ii) shared support of $p_s$ and $p_t$, (iii)
linearity of $T$. Non-linear $T$ (e.g.\ deep mean-teacher) and
heavy-tailed conditionals are deferred to Appendix~D (future work).}

\end{document}
